\chapter{Banach Spaces}

\begin{center}\itshape
In infinite dimension the unit ball is not compact, a linear map can be discontinuous, and a sequence of unit vectors can
converge to zero. What survives of linear algebra, and what new tools replace the ones that break?
\end{center}

On Akademicka Street in Lw\'ow, in the 1930s, there was a caf\'e with marble tables. The Scottish Caf\'e had nothing Scottish
about it except the name, but it had a peculiar clientele: Stefan Banach, Stanis\l aw Mazur, Stanis\l aw Ulam, Hugo Steinhaus
and their friends sat there for hours, sometimes for seventeen hours, and wrote mathematics on the tables in pencil. The waiters
wiped the marble every night, and the theorems were lost. In 1935 Banach's wife, \L ucja, put an end to the waste: she bought a
thick notebook and left it with the head waiter, who brought it to the table on request. The \emph{Scottish Book} collected
almost two hundred problems, each with a prize offered by its author: a small beer, a bottle of wine, a fondue in Geneva.

Problem~153, written by Mazur on 6 November 1936, asked a question about approximating continuous functions of two variables,
which turned out to be equivalent to this: can every compact operator on a Banach space be approximated by operators of finite
rank? The prize was a live goose. The question waited thirty-six years. In 1972 Per Enflo, a young Swede who also gave concerts as
a pianist, constructed a Banach space where the answer is no, travelled to Warsaw, and received the goose from Mazur's own hands,
in a ceremony broadcast on Polish television.

The spaces of that notebook are the subject of this chapter. Banach's book \emph{Th\'eorie des op\'erations lin\'eaires} (1932)
set out their theory: complete normed spaces, which Banach called spaces of ``type~(B)'' and everyone else calls Banach spaces.
Three of its pillars rest on Baire's theorem from Volume~6: a family of operators bounded at each point is bounded uniformly; a
continuous linear surjection is open; a linear map with a closed graph is continuous. A fourth pillar, the Hahn--Banach theorem,
found independently by Hans Hahn in Vienna in 1927 (two years before he supervised a thesis by a student called Kurt G\"odel) and
by Banach in 1929, says that there are enough continuous linear forms to see every vector. With them, this chapter identifies
the duals of the $L^p$ spaces and of $C(K)$, shows that the Fourier series of a continuous function can diverge, and builds the
weak topologies in which bounded sets become compact again.

Throughout, the scalars are $\mathbb K=\R$ or $\C$, and the results hold for both unless said otherwise.

\section{Normed spaces and bounded operators}

A \emph{norm} on a vector space $E$ over $\mathbb K$ is a map $\|\cdot\|:E\to[0,\infty[$ with $\|x\|=0\iff x=0$, $\|\lambda x\|=|\lambda|\|x\|$ and $\|x+y\|\le\|x\|+\|y\|$. It defines the distance
$d(x,y)=\|x-y\|$, and a \emph{Banach space} is a normed space that is complete for it. Volume~6 met many: $\R^n$ with any norm, $C_b(X)$ with the sup norm (Chapter~2), and $L^p(\mu)$ for $1\le p\le\infty$
(Theorem~9.4 there). Three more families will serve as examples throughout:
\begin{itemize}
\item $\ell^p$ for $1\le p<\infty$, the sequences $(x_n)_{n\ge1}$ with $\|x\|_p=(\sum_n|x_n|^p)^{1/p}<\infty$, which is $L^p$ of the counting measure on $\N^*$; and $\ell^\infty$, the bounded sequences with the sup norm;
\item $c_0\subseteq\ell^\infty$, the sequences tending to $0$, a closed subspace, hence a Banach space;
\item $C(K)$ for a compact metric space $K$, with $\|f\|_\infty=\max_K|f|$.
\end{itemize}

A linear map $T:E\to F$ between normed spaces is called an \emph{operator}; it is \emph{bounded} if $\|Tx\|\le C\|x\|$ for some constant $C$. The smallest such $C$ is the \emph{operator norm}
\[\|T\|=\sup_{\|x\|\le1}\|Tx\|.\]
When $F=\mathbb K$, an operator is a \emph{linear form}, and the bounded linear forms on $E$ form the \emph{dual} $E^*$.

\begin{proposition}\label{prop:boundedcont}
For an operator $T:E\to F$, the following are equivalent: (a) $T$ is continuous; (b) $T$ is continuous at $0$; (c) $T$ is bounded. The bounded operators form a normed space $\mathcal L(E,F)$, which is a
Banach space when $F$ is. In particular $E^*$ is always a Banach space.
\end{proposition}
\begin{proof}
(a)$\Rightarrow$(b) is clear. (b)$\Rightarrow$(c): there is $\delta>0$ with $\|Tx\|\le1$ when $\|x\|\le\delta$; for $x\ne0$, apply this to $\delta x/\|x\|$ to get $\|Tx\|\le\frac1\delta\|x\|$. (c)$\Rightarrow$(a):
$\|Tx-Ty\|=\|T(x-y)\|\le C\|x-y\|$, so $T$ is Lipschitz.

The operator norm is a norm (for the triangle inequality, $\|(S+T)x\|\le\|Sx\|+\|Tx\|\le(\|S\|+\|T\|)\|x\|$). Let $F$ be complete and $(T_n)$ Cauchy in $\mathcal L(E,F)$. For each $x$,
$\|T_nx-T_mx\|\le\|T_n-T_m\|\|x\|$, so $(T_nx)$ is Cauchy in $F$; call its limit $Tx$. The map $T$ is linear as a pointwise limit of linear maps. Given $\varepsilon>0$, choose $N$ with $\|T_n-T_m\|\le\varepsilon$ for
$n,m\ge N$; then $\|T_nx-T_mx\|\le\varepsilon\|x\|$, and letting $m\to\infty$, $\|T_nx-Tx\|\le\varepsilon\|x\|$. So $T_n-T$ is bounded, hence so is $T$, and $\|T_n-T\|\le\varepsilon$ for $n\ge N$.
\end{proof}

In finite dimension every operator is bounded, and every norm is equivalent to every other (Volume~6, Corollary~4.6). The next theorem is the first sign that infinite dimension is different, and it explains
why compactness arguments, so effective in $\R^n$, need new ideas in function spaces.

\begin{lemma}[Riesz's lemma]\label{lem:riesz}
Let $Y$ be a closed subspace of a normed space $E$, with $Y\ne E$, and let $0<\theta<1$. There is $x\in E$ with $\|x\|=1$ and $d(x,Y)\ge\theta$.
\end{lemma}
\begin{proof}
Take $z\notin Y$. Since $Y$ is closed, $d=d(z,Y)>0$, and there is $y_0\in Y$ with $\|z-y_0\|\le d/\theta$. Let $x=\frac{z-y_0}{\|z-y_0\|}$. For $y\in Y$, $y_0+\|z-y_0\|y\in Y$, so
\[\|x-y\|=\frac{\|z-(y_0+\|z-y_0\|y)\|}{\|z-y_0\|}\ge\frac d{\|z-y_0\|}\ge\theta.\]
\end{proof}

\begin{theorem}[Riesz]\label{thm:rieszcompact}
The closed unit ball of a normed space $E$ is compact if and only if $E$ has finite dimension.
\end{theorem}
\begin{proof}
If $\dim E=n$, a linear bijection $\mathbb K^n\to E$ is a homeomorphism (Volume~6, Corollary~4.6), and the closed unit ball is closed and bounded, hence compact. If $\dim E=\infty$, build $x_1,x_2,\dots$ inductively:
$Y_n=\operatorname{span}(x_1,\dots,x_n)$ is finite-dimensional, hence complete, hence closed, and $Y_n\ne E$; Lemma~\ref{lem:riesz} gives a unit vector $x_{n+1}$ with $d(x_{n+1},Y_n)\ge\frac12$. Then $\|x_n-x_m\|\ge\frac12$
for $n\ne m$, so $(x_n)$ has no convergent subsequence, and the ball is not compact (Volume~6, Theorem~4.8).
\end{proof}

In $\ell^2$ there is no need for the lemma: the unit vectors $e_n$ satisfy $\|e_n-e_m\|_2=\sqrt2$. The ball is a very big place.

\begin{think}
Is every linear map between Banach spaces continuous? In finite dimension, yes. In infinite dimension, try to imagine a linear form on $\ell^2$ that is not continuous. Can you write one down explicitly?
\end{think}

\section{The Hahn--Banach theorem}

In $\R^n$ there are plenty of linear forms: the coordinates. In a general normed space, it is not even clear that there is a single non-zero bounded linear form. The Hahn--Banach theorem produces them,
extending a form from a subspace without increasing its size. A map $p:E\to\R$ on a real vector space is \emph{sublinear} if $p(x+y)\le p(x)+p(y)$ and $p(tx)=tp(x)$ for $t\ge0$; norms are sublinear, and so
are $x\mapsto\limsup x_n$ on $\ell^\infty$ and the gauges of convex sets below.

\begin{theorem}[Hahn--Banach, analytic form]\label{thm:hb}
Let $E$ be a real vector space, $p$ sublinear on $E$, $G$ a subspace and $g:G\to\R$ linear with $g\le p$ on $G$. There is a linear $f:E\to\R$ extending $g$ with $f\le p$ on $E$.
\end{theorem}
\begin{proof}
\emph{One more dimension.} Let $x_0\notin G$. Extending $g$ to $G\oplus\R x_0$ amounts to choosing $c=f(x_0)$, and the requirement is $g(y)+tc\le p(y+tx_0)$ for all $y\in G$ and $t\in\R$. For $t>0$, dividing by $t$, it
reads $c\le p(y+x_0)-g(y)$ for all $y\in G$; for $t<0$, dividing by $-t$, it reads $c\ge g(y')-p(y'-x_0)$ for all $y'\in G$. A suitable $c$ exists because every lower bound is below every upper bound:
\[g(y)+g(y')=g(y+y')\le p(y+y')\le p(y+x_0)+p(y'-x_0).\]
\emph{All dimensions at once.} Order the pairs $(H,h)$, with $H\supseteq G$ a subspace and $h$ a linear extension of $g$ to $H$ with $h\le p$, by extension. Every chain has an upper bound (the union of the domains,
with the common extension), so Zorn's lemma (Volume~6, Theorem~1.19) gives a maximal pair $(H,h)$. If $H\ne E$, the first step extends $h$ further, contradicting maximality. So $H=E$.
\end{proof}

\begin{remark}
Zorn is not needed when $E$ is a separable normed space and $p\le C\|\cdot\|$: extend one dimension at a time along a dense sequence $(x_n)$, which defines $f$ on the dense subspace $\bigcup_n(G+\operatorname{span}(x_1,\dots,x_n))$, where
$|f|\le C\|\cdot\|$; then $f$ is uniformly continuous and extends by continuity to $E$ (Volume~6, Lemma~2.8), still with $f\le p$, since $p$ is continuous. In general some form of choice is needed: there are models of set
theory without choice in which every linear form on $\ell^\infty$ that vanishes on $c_0$ is zero.
\end{remark}

For complex spaces, a linear form $f$ is determined by its real part: $\operatorname{Im}f(x)=-\operatorname{Re}f(ix)$, since $f(ix)=if(x)$.

\begin{corollary}\label{cor:hbnorm}
Let $E$ be a normed space over $\mathbb K$ and $G$ a subspace. (a) Every $g\in G^*$ extends to $f\in E^*$ with $\|f\|=\|g\|$. (b) For every $x_0\in E$ there is $f\in E^*$ with $\|f\|=1$ and $f(x_0)=\|x_0\|$; hence
$\|x\|=\max_{\|f\|\le1}|f(x)|$, and $E^*$ separates the points of $E$. (c) If $G$ is closed and $x_0\notin G$, there is $f\in E^*$ with $f=0$ on $G$ and $f(x_0)=1$.
\end{corollary}
\begin{proof}
(a) Real case: apply Theorem~\ref{thm:hb} with $p=\|g\|\,\|\cdot\|$; then $f\le p$ and $-f(x)=f(-x)\le p(x)$, so $|f|\le\|g\|\|\cdot\|$. Complex case: extend $u=\operatorname{Re}g$, a real linear form with $|u|\le\|g\|\|\cdot\|$, to a real
linear $U$ on $E$ with the same bound, and set $f(x)=U(x)-iU(ix)$. It is complex linear ($f(ix)=U(ix)-iU(-x)=i(U(x)-iU(ix))=if(x)$) and extends $g$. Given $x$, write $f(x)=e^{i\theta}|f(x)|$; then
$|f(x)|=f(e^{-i\theta}x)=U(e^{-i\theta}x)\le\|g\|\|x\|$, the middle equality because the number is real. (b) For $x_0\ne0$ apply (a) to $g(\lambda x_0)=\lambda\|x_0\|$ on $\mathbb Kx_0$, of norm $1$. Then
$|f(x)|\le\|x\|$ for $\|f\|\le1$, with equality for the $f$ just built. (c) On $G\oplus\mathbb Kx_0$ let $g(y+\lambda x_0)=\lambda$. It is bounded: $|\lambda|\,d(x_0,G)\le\|\lambda x_0+y\|$ for all $y\in G$, and
$d(x_0,G)>0$ since $G$ is closed. Extend by (a).
\end{proof}

The geometric form of the theorem separates convex sets by hyperplanes. Let $C$ be an open convex set containing $0$. Its \emph{gauge} is $p_C(x)=\inf\{t>0:x\in tC\}$. If $B(0,r)\subseteq C$, then $x\in tC$ as soon as
$t>\|x\|/r$, so $p_C(x)\le\|x\|/r$. The gauge is sublinear: positive homogeneity is clear, and if $x\in sC$, $y\in tC$, then $x+y=(s+t)\bigl(\frac s{s+t}\frac xs+\frac t{s+t}\frac yt\bigr)\in(s+t)C$ by convexity. Finally
$C=\{p_C<1\}$: if $x\in C$, then $(1+\varepsilon)x\in C$ for small $\varepsilon$ since $C$ is open, so $p_C(x)\le\frac1{1+\varepsilon}<1$; conversely if $p_C(x)<1$, then $x\in tC$ for some $t<1$, and
$x=t\frac xt+(1-t)0\in C$.

\begin{theorem}[Hahn--Banach, geometric form]\label{thm:separation}
Let $E$ be a real normed space and $A,B$ disjoint non-empty convex sets.
\begin{enumerate}
\item[(a)] If $A$ is open, there are $f\in E^*$ and $\alpha\in\R$ with $f(a)<\alpha\le f(b)$ for all $a\in A$, $b\in B$.
\item[(b)] If $A$ is closed and $B$ is compact, there are $f\in E^*$ and $\alpha<\beta$ with $f(a)\le\alpha<\beta\le f(b)$ for all $a\in A$, $b\in B$.
\end{enumerate}
\end{theorem}
\begin{proof}
(a) First let $B=\{x_0\}$ and $0\in A$. Define $g(tx_0)=t$ on $\R x_0$. For $t>0$, $p_A(tx_0)=tp_A(x_0)\ge t$ because $x_0\notin A=\{p_A<1\}$; for $t\le0$, $g(tx_0)\le0\le p_A(tx_0)$. So $g\le p_A$, and
Theorem~\ref{thm:hb} extends $g$ to $f\le p_A$. Then $f$ is bounded, since $f(x)\le p_A(x)\le\|x\|/r$ and the same for $-x$; and $f(a)\le p_A(a)<1=f(x_0)$ on $A$. In general, $C=A-B+x_0$, with $x_0=b_0-a_0$ for some $b_0\in B$,
$a_0\in A$, is open (a union of translates of $A$), convex, contains $0$, and does not contain $x_0$ because $A\cap B=\varnothing$. The first case gives $f$ with $f(a)-f(b)+f(x_0)<f(x_0)$, that is $f(a)<f(b)$, for all $a,b$. As
$f(A)$ is an open interval (a non-zero linear form is an open map: $f(a+tu)=f(a)+tf(u)$ with $f(u)\ne0$), $\alpha=\sup f(A)$ is not attained, and $f(a)<\alpha\le f(b)$.

(b) The function $b\mapsto d(b,A)$ is continuous and positive on the compact $B$ (as $A$ is closed), so it has a minimum $2r>0$. Then $A_r=A+B(0,r)$ and $B_r=B+B(0,r)$ are open, convex and disjoint. By (a) there is $f\ne0$
with $f(a+u)\le f(b+v)$ for all $a\in A$, $b\in B$, $\|u\|,\|v\|<r$. Taking the supremum over $u$ and the infimum over $v$, $f(a)+r\|f\|\le f(b)-r\|f\|$; take $\alpha=\sup_Af$ and $\beta=\alpha+2r\|f\|$.
\end{proof}

\begin{figure}[ht]
\centering
\fig{ch01-separation}
\caption{The geometric Hahn--Banach theorem: a closed convex set and a compact convex set that do not meet are separated by a closed hyperplane $\{f=\gamma\}$, with a strip of positive width between them.}
\label{fig:separation}
\end{figure}

\begin{beauty}
A closed convex set is the intersection of the closed half-spaces that contain it: if $x_0\notin C$, part (b) with $B=\{x_0\}$ gives a half-space containing $C$ and not $x_0$. A convex body, described by its points, is
equally well described by its supporting hyperplanes, that is, by linear forms. This duality is the heart of convex optimisation, where minimisation problems are paired with dual maximisation problems.
\end{beauty}

\section{Three consequences of Baire's theorem}

Baire's theorem (Volume~6, Theorem~5.2) says that a complete metric space is not a countable union of closed sets with empty interior. Banach and Steinhaus saw in 1927 that, applied to the right sets, it turns pointwise
information about operators into uniform information.

\begin{theorem}[uniform boundedness, Banach--Steinhaus]\label{thm:ubp}
Let $E$ be a Banach space, $F$ a normed space and $(T_i)_{i\in I}$ a family in $\mathcal L(E,F)$.
\begin{enumerate}
\item[(a)] If $\sup_i\|T_ix\|<\infty$ for every $x\in E$, then $\sup_i\|T_i\|<\infty$.
\item[(b)] If $\sup_i\|T_i\|=\infty$, then the set of $x$ with $\sup_i\|T_ix\|=\infty$ is a dense $G_\delta$ of $E$.
\end{enumerate}
\end{theorem}
\begin{proof}
Let $F_n=\{x:\|T_ix\|\le n\text{ for all }i\}$, closed as an intersection of closed sets. If some $F_n$ contains a ball $\overline B(x_0,r)$, then for $\|z\|\le r$ and every $i$,
$\|T_iz\|\le\|T_i(x_0+z)\|+\|T_ix_0\|\le2n$, so $\|T_i\|\le\frac{2n}r$ for all $i$. Hence, if $\sup_i\|T_i\|=\infty$, every $F_n$ has empty interior, and $\bigcup_nF_n$ is a countable union of closed sets with empty
interior; by Baire's theorem its complement $\{x:\sup_i\|T_ix\|=\infty\}=\bigcap_n(E\setminus F_n)$, a countable intersection of dense open sets, is dense. This proves (b), and (a) follows: under its hypothesis
$\bigcup_nF_n=E$, so the norms cannot be unbounded.
\end{proof}

The first application settles a question that occupied the nineteenth century. Dirichlet proved in 1829 that the Fourier series of a piecewise smooth function converges, and many believed that continuity would be
enough. In 1873 Paul du Bois-Reymond built a continuous function whose Fourier series diverges at a point. With Theorem~\ref{thm:ubp}, no construction is needed, and the result is stronger.

\begin{theorem}\label{thm:fourierdiverge}
For $f\in C(\mathbb T)$, the continuous $2\pi$-periodic functions with the sup norm, let $S_Nf(0)=\sum_{|n|\le N}\hat f(n)$ be the partial sums of its Fourier series at $0$. The set of $f$ with
$\sup_N|S_Nf(0)|=\infty$ is a dense $G_\delta$ of $C(\mathbb T)$. In particular, the Fourier series of most continuous functions diverge at $0$.
\end{theorem}
\begin{proof}
By Volume~4, Chapter~6, $\Lambda_Nf=S_Nf(0)=\frac1{2\pi}\int_{-\pi}^\pi f(t)D_N(t)\,dt$ with the Dirichlet kernel $D_N(t)=\frac{\sin((N+\frac12)t)}{\sin(t/2)}$. Clearly $|\Lambda_Nf|\le L_N\|f\|_\infty$ with
$L_N=\frac1{2\pi}\int_{-\pi}^\pi|D_N|$. Conversely, continuous functions $f_k$ with values in $[-1,1]$ converge to $\operatorname{sgn}D_N$ almost everywhere (take $f_k=\max(-1,\min(1,kD_N))$), and dominated
convergence gives $\Lambda_Nf_k\to L_N$; so $\|\Lambda_N\|=L_N$. Using $|\sin(t/2)|\le|t|/2$ and the substitution $u=(N+\frac12)t$,
\[L_N\ge\frac2\pi\int_0^\pi\frac{|\sin((N+\frac12)t)|}t\,dt\ge\frac2\pi\sum_{k=0}^{N-1}\int_{k\pi}^{(k+1)\pi}\frac{|\sin u|}{(k+1)\pi}\,du=\frac4{\pi^2}\sum_{k=1}^N\frac1k,\]
which tends to infinity. Theorem~\ref{thm:ubp}(b) concludes.
\end{proof}

The numbers $L_N$ are the \emph{Lebesgue constants}; they grow like $\frac4{\pi^2}\ln N$, and a direct computation gives them. Kolmogorov went further in 1926 and built an integrable function whose Fourier series diverges
\emph{everywhere}. In the other direction, Lennart Carleson proved in 1966 that the Fourier series of every $f\in L^2$ converges almost everywhere, one of the deepest theorems of analysis. Chapter~2 will show that in
the $L^2$ norm, convergence holds for every $f\in L^2$, and Fej\'er's theorem (Volume~4) that Ces\`aro means converge uniformly for continuous $f$: divergence is a fact about the partial sums, not about the series.

\begin{theorem}[open mapping]\label{thm:openmapping}
Let $E$ and $F$ be Banach spaces and $T\in\mathcal L(E,F)$ surjective. Then $T$ is open: there is $\delta>0$ with $B_F(0,\delta)\subseteq T(B_E(0,1))$. If $T$ is moreover injective, $T^{-1}$ is bounded.
\end{theorem}
\begin{proof}
\emph{Step 1: the closure of $T(B_E(0,1))$ contains a ball.} Since $T$ is onto, $F=\bigcup_n\overline{T(B_E(0,n))}$, and by Baire one of these closed sets, hence (scaling) $\overline{T(B_E(0,1))}$, contains a ball
$B_F(y_0,4\delta)$. As $\overline{T(B_E(0,1))}$ is convex and symmetric, it contains $\frac12(B_F(y_0,4\delta)+B_F(-y_0,4\delta))\supseteq B_F(0,4\delta)$, and by scaling $\overline{T(B_E(0,2^{-k}))}\supseteq B_F(0,2^{2-k}\delta)$ for all $k\ge0$.

\emph{Step 2: remove the closure.} Let $\|y\|<\delta$. By Step~1 with $k=2$ there is $x_1$ with $\|x_1\|<\frac14$ and $\|y-Tx_1\|<\frac\delta2$; then $x_2$ with $\|x_2\|<\frac18$ and $\|y-Tx_1-Tx_2\|<\frac\delta4$; and so on,
$\|x_k\|<2^{-k-1}$, $\|y-T(x_1+\cdots+x_k)\|<2^{-k}\delta$. The series $x=\sum_kx_k$ converges absolutely in the Banach space $E$, with $\|x\|<\frac12<1$, and $Tx=y$ by continuity of $T$. So $B_F(0,\delta)\subseteq T(B_E(0,1))$.

If $T$ is bijective, $\|Tx\|<\delta$ implies $\|x\|<1$, so $\|T^{-1}y\|\le\frac1\delta\|y\|$.
\end{proof}

\begin{theorem}[closed graph]\label{thm:closedgraph}
Let $E$ and $F$ be Banach spaces and $T:E\to F$ linear. If the graph $\{(x,Tx):x\in E\}$ is closed in $E\times F$, then $T$ is bounded.
\end{theorem}
\begin{proof}
With the norm $\|x\|+\|y\|$, $E\times F$ is a Banach space, and the graph $\Gamma$, a closed subspace, is one too. The projection $\pi:\Gamma\to E$, $(x,Tx)\mapsto x$, is bounded and bijective, so $\pi^{-1}$ is bounded by
Theorem~\ref{thm:openmapping}: $\|x\|+\|Tx\|\le C\|x\|$. Hence $\|Tx\|\le C\|x\|$.
\end{proof}

In practice the closed graph theorem says that to prove continuity it suffices to check that if $x_n\to x$ \emph{and} $Tx_n\to y$, then $y=Tx$: one may assume the convergence one wants to prove. The theorem fails without
completeness: differentiation, defined on $C^1[0,1]$ with the sup norm, has a closed graph and is unbounded. And in the models of set theory mentioned after Theorem~\ref{thm:hb}, every linear map between Banach spaces is
continuous (Garnir, Wright): the discontinuous linear forms of Exercise~\ref{ex:discontinuous} cannot be built without a Hamel basis.

\section{Duals of the classical spaces}

To know a Banach space is, to a large extent, to know its dual. Hölder's inequality (Volume~6, Theorem~9.2) shows that every $g\in L^q$ defines a bounded linear form $\varphi_g(f)=\int fg\,d\mu$ on $L^p$, $\frac1p+\frac1q=1$,
with $\|\varphi_g\|\le\|g\|_q$. Frigyes Riesz proved in 1910 that there are no others when $1<p<\infty$; the case $p=1$ is due to Steinhaus (1919) for $[0,1]$. With the Radon--Nikodym theorem the proof is short.

\begin{theorem}[duality of $L^p$]\label{thm:lpdual}
Let $\mu$ be a $\sigma$-finite measure, $1\le p<\infty$ and $q$ the conjugate exponent. The map $g\mapsto\varphi_g$ is an isometric isomorphism from $L^q(\mu)$ onto $L^p(\mu)^*$.
\end{theorem}
\begin{proof}
We treat real scalars first.

\emph{Isometry.} Let $g\in L^q$, $g\ne0$. If $p>1$, put $f=|g|^{q-1}\operatorname{sgn}g$; then $|f|^p=|g|^q$, so $\|f\|_p=\|g\|_q^{q/p}$ and $\varphi_g(f)=\|g\|_q^q=\|g\|_q\|f\|_p$. If $p=1$ and $0<\varepsilon<\|g\|_\infty$, the set
$\{|g|>\|g\|_\infty-\varepsilon\}$ has positive measure and, $\mu$ being $\sigma$-finite, contains a set $A$ with $0<\mu(A)<\infty$; with $f=\mathbf 1_A\operatorname{sgn}g$, $\varphi_g(f)\ge(\|g\|_\infty-\varepsilon)\|f\|_1$. In both
cases $\|\varphi_g\|\ge\|g\|_q$.

\emph{Surjectivity when $\mu$ is finite.} Let $\varphi\in(L^p)^*$ and $\nu(A)=\varphi(\mathbf 1_A)$. If $A=\bigsqcup_kA_k$, then $\|\mathbf 1_A-\mathbf 1_{A_1\cup\cdots\cup A_n}\|_p=\mu(\bigcup_{k>n}A_k)^{1/p}\to0$ because $p<\infty$, so
$\nu(A)=\sum_k\nu(A_k)$: $\nu$ is a finite signed measure, and $\nu(A)=0$ when $\mu(A)=0$. Its Jordan parts $\nu^\pm(A)=\pm\nu(A\cap P^\pm)$ (Volume~6, Corollary~10.3) are finite measures that also vanish on $\mu$-null sets, so the
Radon--Nikodym theorem (Volume~6, Theorem~10.4) gives $g=\frac{d\nu^+}{d\mu}-\frac{d\nu^-}{d\mu}\in L^1(\mu)$ with $\varphi(\mathbf 1_A)=\int_Ag\,d\mu$. By linearity $\varphi(f)=\int fg\,d\mu$ for simple $f$, and then for bounded measurable
$f$, which are uniform limits of simple functions $f_n$ (uniform convergence implies $L^p$ convergence on a finite measure space, and $\int f_ng\to\int fg$ because $g\in L^1$).

Now $g\in L^q$. If $p>1$, let $f_n=|g|^{q-1}\operatorname{sgn}g\,\mathbf 1_{\{|g|\le n\}}$, bounded; then
\[\int_{\{|g|\le n\}}|g|^q=\varphi(f_n)\le\|\varphi\|\,\|f_n\|_p=\|\varphi\|\Bigl(\int_{\{|g|\le n\}}|g|^q\Bigr)^{1/p},\]
so $(\int_{\{|g|\le n\}}|g|^q)^{1/q}\le\|\varphi\|$, and monotone convergence gives $\|g\|_q\le\|\varphi\|$. If $p=1$, let $A=\{g>\|\varphi\|+\varepsilon\}$; then $(\|\varphi\|+\varepsilon)\mu(A)\le\int_Ag=\varphi(\mathbf 1_A)\le\|\varphi\|\mu(A)$,
so $\mu(A)=0$, and likewise for $\{g<-\|\varphi\|-\varepsilon\}$: $\|g\|_\infty\le\|\varphi\|$. Finally $\varphi$ and $\varphi_g$ are bounded and agree on the simple functions, which are dense in $L^p$ (Volume~6, Proposition~9.6), so
$\varphi=\varphi_g$.

\emph{The $\sigma$-finite case.} Let $E_1\subseteq E_2\subseteq\cdots$ have finite measure and union $E$. The finite case on each $E_n$ gives $g_n$ on $E_n$, with $\|g_n\|_q\le\|\varphi\|$, and by uniqueness $g_{n+1}=g_n$ a.e. on $E_n$. The
function $g$ equal to $g_n$ on $E_n$ satisfies $\|g\|_q\le\|\varphi\|$ by monotone convergence, and $\varphi(f)=\int fg$ for every $f$ vanishing outside some $E_n$. Such $f$ are dense in $L^p$ (dominated convergence:
$f\mathbf 1_{E_n}\to f$), so $\varphi=\varphi_g$.

\emph{Complex scalars.} Let $\varphi$ be a bounded complex-linear form on $L^p$. On real-valued functions, both $u\mapsto\operatorname{Re}\varphi(u)$ and $u\mapsto\operatorname{Im}\varphi(u)$ are real-linear and bounded by $\|\varphi\|\,\|u\|_p$. The real case just proved supplies real functions $g_{\mathrm R},g_{\mathrm I}\in L^q$ with $\operatorname{Re}\varphi(u)=\int ug_{\mathrm R}$ and $\operatorname{Im}\varphi(u)=\int ug_{\mathrm I}$. Set $g=g_{\mathrm R}+ig_{\mathrm I}$. For real $u$, $\varphi(u)=\int ug$. For $f=u+iv$, complex linearity gives $\varphi(f)=\varphi(u)+i\varphi(v)=\int fg$. As in the isometry paragraph, $f=|g|^{q-1}\operatorname{sgn}g$ gives $\|g\|_q\le\|\varphi\|$, and H\"older gives the reverse. So the map is an isometric isomorphism over $\C$ as well.
\end{proof}

With the counting measure, $(\ell^p)^*=\ell^q$ for $1\le p<\infty$. The case $p=\infty$ is excluded for a good reason: $(L^\infty)^*$ is much larger than $L^1$ (Exercise~\ref{ex:linftydual}), and its extra elements, like the
Banach limits of Exercise~\ref{ex:banachlimit}, can only be produced with choice.

\begin{proposition}\label{prop:c0dual}
Every $\varphi\in c_0^*$ is $\varphi(x)=\sum_na_nx_n$ for a unique $a\in\ell^1$, and $\|\varphi\|=\|a\|_1$: $c_0^*=\ell^1$.
\end{proposition}
\begin{proof}
Let $a_n=\varphi(e_n)$. For $x=\sum_{n\le N}\overline{\operatorname{sgn}a_n}\,e_n$, of norm at most $1$, $\varphi(x)=\sum_{n\le N}|a_n|\le\|\varphi\|$; so $a\in\ell^1$ with $\|a\|_1\le\|\varphi\|$. Every $x\in c_0$ is the limit in sup norm of its
truncations $\sum_{n\le N}x_ne_n$ (this is where $x_n\to0$ is used), so $\varphi(x)=\lim_N\sum_{n\le N}a_nx_n=\sum_na_nx_n$, and $|\varphi(x)|\le\|a\|_1\|x\|_\infty$.
\end{proof}

For $C(K)$ the dual consists of measures. A \emph{finite signed Borel measure} $\nu$ on $K$ (Volume~6, Chapter~10) has a Jordan decomposition $\nu=\nu^+-\nu^-$ into mutually singular finite measures; its \emph{total variation}
is $\|\nu\|=\nu^+(K)+\nu^-(K)$. Let $M(K)$ be the space of these measures; by the next theorem it is a Banach space for this norm, isometric to $C(K)^*$.

\begin{theorem}[Riesz--Markov]\label{thm:rieszmarkov}
Let $K$ be a compact metric space. Every $\varphi\in C(K,\R)^*$ is $\varphi(f)=\int f\,d\nu$ for a unique finite signed Borel measure $\nu$, and $\|\varphi\|=\|\nu\|$. So $C(K)^*=M(K)$ isometrically.
\end{theorem}
\begin{proof}
\emph{A positive part.} For $f\ge0$ in $C(K)$ let $\varphi^+(f)=\sup\{\varphi(g):g\in C(K),\ 0\le g\le f\}$; then $0\le\varphi^+(f)\le\|\varphi\|\|f\|_\infty$. It is additive on non-negative functions: $\ge$ is clear, since
$0\le g_1\le f_1$ and $0\le g_2\le f_2$ give $0\le g_1+g_2\le f_1+f_2$; for $\le$, if $0\le g\le f_1+f_2$, then $g_1=\min(g,f_1)$ and $g_2=g-g_1=\max(g-f_1,0)$ satisfy $0\le g_1\le f_1$, $0\le g_2\le f_2$, and
$\varphi(g)=\varphi(g_1)+\varphi(g_2)\le\varphi^+(f_1)+\varphi^+(f_2)$. It is positively homogeneous. Writing $f=f^+-f^-$, the formula $\varphi^+(f)=\varphi^+(f^+)-\varphi^+(f^-)$ is well defined (if $f=u-v$ with $u,v\ge0$,
then $u+f^-=v+f^+$ and additivity gives $\varphi^+(u)-\varphi^+(v)=\varphi^+(f^+)-\varphi^+(f^-)$) and linear. So $\varphi^+$ is a positive linear form, and so is $\varphi^-=\varphi^+-\varphi$, since $\varphi^+(f)\ge\varphi(f)$ for $f\ge0$
(take $g=f$).

\emph{Existence.} The Riesz representation theorem (Volume~6, Theorem~11.3) gives finite measures $\mu^\pm$ with $\varphi^\pm(f)=\int f\,d\mu^\pm$; then $\nu=\mu^+-\mu^-$ represents $\varphi$. \emph{Uniqueness.} If $\int f\,d\nu=0$ for all
$f$, with Jordan decomposition $\nu=\nu^+-\nu^-$, then the finite measures $\nu^+$ and $\nu^-$ have the same integrals, so they are equal by the uniqueness in Riesz's theorem, and $\nu=0$.

\emph{The norm.} $|\int f\,d\nu|\le\int|f|\,d\nu^++\int|f|\,d\nu^-\le\|f\|_\infty\|\nu\|$. Conversely let $K=P\sqcup N$ be a Hahn decomposition for $\nu$ (Volume~6, Theorem~10.2), $\lambda=\nu^++\nu^-$ and $\varepsilon>0$. Finite Borel
measures on a metric space are regular (Volume~6, Lemma~10.10), so there are closed sets $F_P\subseteq P$, $F_N\subseteq N$ with $\lambda(P\setminus F_P)<\varepsilon$, $\lambda(N\setminus F_N)<\varepsilon$. The function
$h(x)=\frac{d(x,F_N)-d(x,F_P)}{d(x,F_N)+d(x,F_P)}$ is continuous (the closed sets $F_P,F_N$ are disjoint, so the denominator does not vanish), with values in $[-1,1]$, equal to $1$ on $F_P$ and $-1$ on $F_N$. So
\[\varphi(h)=\int h\,d\nu^+-\int h\,d\nu^-\ge(\nu^+(F_P)-\varepsilon)-(-\nu^-(F_N)+\varepsilon)\ge\|\nu\|-4\varepsilon,\]
using $\nu^+(K\setminus F_P)\le\lambda(P\setminus F_P)+\nu^+(N)<\varepsilon$ (as $\nu^+(N)=0$) and similarly for $\nu^-$. Hence $\|\varphi\|\ge\|\nu\|$.
\end{proof}

A Banach space is \emph{reflexive} if the canonical map $J:E\to E^{**}$, $J(x)(\varphi)=\varphi(x)$, is onto. By Corollary~\ref{cor:hbnorm}(b), $J$ is always an isometry: $\|J(x)\|=\sup_{\|\varphi\|\le1}|\varphi(x)|=\|x\|$. So $E$ sits inside
its bidual, and the question is whether it fills it.

\begin{proposition}\label{prop:reflexive}
$L^p(\mu)$ is reflexive for $1<p<\infty$ and $\sigma$-finite $\mu$. The space $c_0$ is not reflexive.
\end{proposition}
\begin{proof}
Let $\Phi\in(L^p)^{**}$. Then $g\mapsto\Phi(\varphi_g)$ is a bounded linear form on $L^q$, so by Theorem~\ref{thm:lpdual} (with $q<\infty$) there is $f\in L^p$ with $\Phi(\varphi_g)=\int fg=\varphi_g(f)$ for all $g$. Since every element of
$(L^p)^*$ is some $\varphi_g$, $\Phi=J(f)$. For $c_0$: by Proposition~\ref{prop:c0dual} and Theorem~\ref{thm:lpdual}, $c_0^{**}=(\ell^1)^*=\ell^\infty$, and under these identifications $J$ is the inclusion $c_0\subseteq\ell^\infty$, which
misses the constant sequence $1$.
\end{proof}

\section{Weak topologies}

The norm topology of an infinite-dimensional space has too many open sets for compactness to be common (Theorem~\ref{thm:rieszcompact}). The remedy is to weaken the topology. A sequence $x_n$ \emph{converges weakly} to $x$,
written $x_n\rightharpoonup x$, if $\varphi(x_n)\to\varphi(x)$ for every $\varphi\in E^*$; a sequence $\varphi_n$ in $E^*$ \emph{converges weak*} to $\varphi$, written $\varphi_n\overset*\rightharpoonup\varphi$, if $\varphi_n(x)\to\varphi(x)$ for every
$x\in E$. The \emph{weak*} topology on $E^*$ is the coarsest topology making all the maps $\varphi\mapsto\varphi(x)$ continuous, that is, the topology induced by the product topology of $\mathbb K^E$ (Volume~6, Chapter~3).

\begin{proposition}\label{prop:weakbounded}
A weakly convergent sequence is bounded, and if $x_n\rightharpoonup x$ then $\|x\|\le\liminf\|x_n\|$. If $E$ is a Banach space, a weak* convergent sequence in $E^*$ is bounded.
\end{proposition}
\begin{proof}
If $x_n\rightharpoonup x$, the forms $J(x_n)$ on the Banach space $E^*$ (Proposition~\ref{prop:boundedcont}) are bounded at every point $\varphi$, so $\sup_n\|x_n\|=\sup_n\|J(x_n)\|<\infty$ by Theorem~\ref{thm:ubp}. With $\varphi$ as in
Corollary~\ref{cor:hbnorm}(b), $\|x\|=\varphi(x)=\lim\varphi(x_n)\le\liminf\|x_n\|$. The weak* statement is Theorem~\ref{thm:ubp} on $E$ itself.
\end{proof}

In $\ell^p$ with $1<p<\infty$, the unit vectors satisfy $e_n\rightharpoonup0$: by Theorem~\ref{thm:lpdual}, every $\varphi$ is $\varphi(x)=\sum_ka_kx_k$ with $a\in\ell^q$, and $\varphi(e_n)=a_n\to0$. Yet $\|e_n\|=1$: the norm can drop
in the limit. In function spaces, weak convergence is the language of oscillation. If $f_n(x)=\operatorname{sgn}\sin(nx)$ on $[0,2\pi]$, then $\|f_n\|_\infty=1$, but $f_n\overset*\rightharpoonup0$ in $L^\infty=(L^1)^*$: for an interval
$I$, $\int_If_n$ is the integral over at most one incomplete half-period plus complete periods that cancel, so $|\int_If_n|\le\frac\pi n$; by linearity $\int f_ng\to0$ for step functions $g$, and then for all $g\in L^1$ by
density (Volume~6, Theorem~9.7) and $|\int f_n(g-h)|\le\|g-h\|_1$. The functions oscillate faster and faster, and their averages against any integrable weight vanish (Figure~\ref{fig:oscillation}).

\begin{figure}[ht]
\centering
\fig{ch01-oscillation}
\caption{$f_n=\operatorname{sgn}\sin(nx)$ for $n=2$ and $n=12$. Each has sup norm $1$, but against any integrable weight $g$ the integral $\int f_ng$ tends to $0$: $f_n\overset*\rightharpoonup0$ in $L^\infty$. In the weak*
limit, the oscillations average out.}
\label{fig:oscillation}
\end{figure}

\begin{theorem}[Banach--Alaoglu]\label{thm:alaoglu}
The closed unit ball $B^*$ of $E^*$ is compact for the weak* topology. If $E$ is separable, every bounded sequence in $E^*$ has a weak* convergent subsequence.
\end{theorem}
\begin{proof}
Each $\varphi\in B^*$ satisfies $|\varphi(x)|\le\|x\|$, so $B^*$ is a subset of $P=\prod_{x\in E}D_x$, where $D_x=\{\lambda\in\mathbb K:|\lambda|\le\|x\|\}$ is compact; and the weak* topology on $B^*$ is the topology induced by
the product. By Tychonoff's theorem (Volume~6, Theorem~4.11), $P$ is compact, so it suffices that $B^*$ is closed in $P$. It is the set of $\omega\in P$ with $\omega(x+y)=\omega(x)+\omega(y)$ and $\omega(\lambda x)=\lambda\omega(x)$ for all
$x,y,\lambda$: an intersection of closed sets, since each condition involves finitely many coordinates, which are continuous functions on $P$. (Such $\omega$ are linear with $|\omega(x)|\le\|x\|$, so in $B^*$.)

If $E$ is separable, let $(x_k)$ be dense and $(\varphi_n)$ a sequence with $\|\varphi_n\|\le M$. By diagonal extraction (as in Volume~6, Theorem~4.15) a subsequence $\psi_n$ has $\psi_n(x_k)$ convergent for every $k$. For $x\in E$ and
$\varepsilon>0$ pick $x_k$ with $\|x-x_k\|<\varepsilon$; then $|\psi_n(x)-\psi_m(x)|\le2M\varepsilon+|\psi_n(x_k)-\psi_m(x_k)|$, so $(\psi_n(x))$ is Cauchy. Its limit $\psi(x)$ is linear in $x$ with $|\psi(x)|\le M\|x\|$, and
$\psi_n\overset*\rightharpoonup\psi$.
\end{proof}

For $K$ compact metric, $C(K)$ is separable (Volume~6, Lemma~11.6) and $C(K)^*=M(K)$, so a sequence of probability measures on $K$ has a subsequence that converges weak*, that is, weakly in the sense of probability; the limit
is positive with mass $1$ (test against $f\ge0$ and $f=1$). This is Prokhorov's theorem on a compact space (Volume~6, Theorem~11.7), now a corollary of a general principle: \emph{bounded sets of a dual space are weak* compact}.
In reflexive spaces the same holds for bounded sets and the weak topology (Exercise~\ref{ex:reflexiveweak}), which is how the calculus of variations finds minimisers and how Chapter~19 finds weak solutions of differential
equations.

\begin{think}
Can a sequence converge weakly but not in norm? In $\ell^2$, yes. In $\ell^1$, surprisingly, no: Schur proved in 1921 that a weakly convergent sequence in $\ell^1$ converges in norm. Why does the unit vector example fail in
$\ell^1$? (Which bounded sequence $a\in\ell^\infty$ tests $e_n$?)
\end{think}

\section{Exercises}

\begin{exercise}[$\star\star$]\label{ex:banachlimit}
\emph{Banach limits.} On real $\ell^\infty$ let $p(x)=\limsup_n\frac{x_1+\cdots+x_n}n$ and let $S$ be the shift, $(Sx)_n=x_{n+1}$. (a) Show that $p$ is sublinear and that there is a linear form $L$ on $\ell^\infty$ with $L\le p$ and
$L(x)=\lim_nx_n$ for every convergent $x$. (b) Show that $L(Sx)=L(x)$, that $L(x)\ge0$ when $x\ge0$, and that $\|L\|=1$. (c) Compute $L(1,0,1,0,\dots)$. (d) Show that $L$ is not of the form $x\mapsto\sum_na_nx_n$ with
$a\in\ell^1$; so $(\ell^\infty)^*\ne\ell^1$.
\end{exercise}

\begin{exercise}[$\star$]\label{ex:discontinuous}
Let $E$ be an infinite-dimensional normed space. Using a Hamel basis (Volume~6, Theorem~1.21), build a linear form on $E$ that is not continuous. Show that its kernel is dense in $E$.
\end{exercise}

\begin{exercise}[$\star$]\label{ex:twonorms}
(a) Let $E$ be a Banach space for two norms with $\|x\|_a\le C\|x\|_b$. Show that the norms are equivalent. (b) Deduce that $C[0,1]$ is not complete for $\|f\|_1=\int_0^1|f|$.
\end{exercise}

\begin{exercise}[$\star\star$]\label{ex:linftydual}
Show that there is $\Lambda\in L^\infty([0,1])^*$ with $\Lambda(f)=f(0)$ for every continuous $f$, and that $\Lambda$ is not of the form $f\mapsto\int fg$ with $g\in L^1$. So $L^1$ is not reflexive: its bidual $(L^\infty)^*$ is
strictly larger than $L^1$.
\end{exercise}

\begin{exercise}[$\star\star$]\label{ex:mazur}
\emph{Mazur's lemma.} Let $x_n\rightharpoonup x$ in a normed space. Show that $x$ is a norm limit of convex combinations of the $x_n$. Illustrate with $e_n\rightharpoonup0$ in $\ell^2$.
\end{exercise}

\begin{exercise}[$\star\star$]\label{ex:reflexiveweak}
Let $E$ be a reflexive Banach space whose dual $E^*$ is separable. Show that every bounded sequence in $E$ has a weakly convergent subsequence. Apply this to $L^p(\R^d)$, $1<p<\infty$.
\end{exercise}

\begin{exercise}[$\star\star\star$]\label{ex:l1quotient}
\emph{Every separable Banach space is a quotient of $\ell^1$.} Let $E$ be a separable Banach space and $(x_n)$ dense in its closed unit ball. Show that $T(a)=\sum_na_nx_n$ defines $T\in\mathcal L(\ell^1,E)$ with $\|T\|\le1$, and that
the closed unit ball of $E$ is contained in the image of the ball of radius $2$ of $\ell^1$. So $T$ is onto and open, and $E$ is isomorphic to $\ell^1/\ker T$.
\end{exercise}

\section{Hints}

\begin{hint}{ex:banachlimit}
(a) Ces\`aro means of a convergent sequence converge to its limit. (b) $\frac1n\sum_{k\le n}(x_k-x_{k+1})=\frac{x_1-x_{n+1}}n$. (c) Add the sequence to its shift. (d) What is $L(e_n)$?
\end{hint}
\begin{hint}{ex:discontinuous}
Normalise the basis vectors and send countably many of them to $1,2,3,\dots$. A hyperplane is either closed or dense, and a linear form with closed kernel is continuous.
\end{hint}
\begin{hint}{ex:twonorms}
(a) The identity from $(E,\|\cdot\|_b)$ to $(E,\|\cdot\|_a)$ is a bounded bijection. (b) Compare with the sup norm on $x^n$.
\end{hint}
\begin{hint}{ex:linftydual}
Continuous functions sit isometrically in $L^\infty$; extend $f\mapsto f(0)$ by Hahn--Banach. Test against continuous bumps concentrating at $0$.
\end{hint}
\begin{hint}{ex:mazur}
The closed convex hull of $\{x_n\}$ is closed and convex; separate $x$ from it if it is outside.
\end{hint}
\begin{hint}{ex:reflexiveweak}
Apply the sequential Banach--Alaoglu theorem to $J(x_n)$ in $E^{**}=(E^*)^*$.
\end{hint}
\begin{hint}{ex:l1quotient}
Given $\|y\|\le1$, approximate $y$ by some $x_{n_1}$ within $\frac12$, then $2(y-x_{n_1})$ by some $x_{n_2}$ within $\frac12$, and so on.
\end{hint}

\section{Solutions}

\begin{solution}{ex:banachlimit}
\idea A ``limit'' for every bounded sequence must be invented, and Hahn--Banach invents it, under a sublinear bound chosen so that the properties wanted are forced. The Ces\`aro $\limsup$ is shift-invariant up to $o(1)$, which is
exactly what is needed.
\writeup (a) $\limsup(a_n+b_n)\le\limsup a_n+\limsup b_n$ and $\limsup ta_n=t\limsup a_n$ for $t\ge0$, so $p$ is sublinear. On the subspace $c$ of convergent sequences let $g(x)=\lim x_n$; Ces\`aro means converge to the same
limit (Volume~0), so $g=p$ on $c$, and Theorem~\ref{thm:hb} gives $L$. (b) For $y=x-Sx$, $\frac1n\sum_{k\le n}y_k=\frac{x_1-x_{n+1}}n\to0$, so $p(y)=p(-y)=0$, hence $L(y)\le0$ and $-L(y)=L(-y)\le0$: $L(y)=0$. If $x\ge0$,
$-L(x)=L(-x)\le p(-x)\le0$. Finally $L(x)\le p(x)\le\sup_nx_n\le\|x\|_\infty$, and $-L(x)\le\|x\|_\infty$ likewise; with $L(1,1,\dots)=1$, $\|L\|=1$. (c) If $x=(1,0,1,0,\dots)$, then $x+Sx=(1,1,\dots)$, so $2L(x)=1$
and $L(x)=\frac12$. (d) Each $e_n$ converges to $0$, so $L(e_n)=0$. If $L(x)=\sum a_nx_n$, then $a_n=L(e_n)=0$ for all $n$ and $L=0$, but $L(1,1,\dots)=1$.
\end{solution}

\begin{solution}{ex:discontinuous}
\idea A linear form is free on a basis, and nothing forces it to stay bounded on unit vectors. Continuity is equivalent to a closed kernel, so an unbounded form has a kernel that is everywhere.
\writeup Let $(e_i)_{i\in I}$ be a Hamel basis with $\|e_i\|=1$; $I$ is infinite. Choose distinct $i_1,i_2,\dots$ and define $f(e_{i_n})=n$, $f(e_i)=0$ for the other $i$, extended linearly. Then $|f(e_{i_n})|=n$ while
$\|e_{i_n}\|=1$, so $f$ is unbounded, hence not continuous. Its kernel $H$ is a hyperplane: $E=H\oplus\mathbb Kx_0$ for any $x_0$ with $f(x_0)=1$. The closure $\overline H$ is a subspace containing $H$, so $\overline H=H$
or $\overline H=E$. If $H$ were closed, then $d=d(x_0,H)>0$ and $|f(h+\lambda x_0)|=|\lambda|\le\frac1d\|h+\lambda x_0\|$ (as in Corollary~\ref{cor:hbnorm}(c)), so $f$ would be bounded. Hence $\overline H=E$.
\end{solution}

\begin{solution}{ex:twonorms}
\idea Comparable complete norms cannot be inequivalent: the open mapping theorem turns the inequality around. So to show a norm is not complete, compare it with a complete norm it does not match.
\writeup (a) The identity $(E,\|\cdot\|_b)\to(E,\|\cdot\|_a)$ is bounded, linear and bijective between Banach spaces, so its inverse is bounded (Theorem~\ref{thm:openmapping}): $\|x\|_b\le C'\|x\|_a$. (b) $\|f\|_1\le\|f\|_\infty$ and
$C[0,1]$ is complete for $\|\cdot\|_\infty$. If it were complete for $\|\cdot\|_1$, the norms would be equivalent by (a); but $f_n(x)=x^n$ has $\|f_n\|_\infty=1$ and $\|f_n\|_1=\frac1{n+1}$.
\end{solution}

\begin{solution}{ex:linftydual}
\idea The evaluation at a point makes sense on continuous functions, not on classes of functions; Hahn--Banach extends it anyway, and the extension cannot be an integral because integrals do not see single points.
\writeup For $f$ continuous on $[0,1]$, the essential supremum of $|f|$ is its maximum (if $|f(x_0)|>\|f\|_\infty+\varepsilon$, continuity gives an interval of positive measure where $|f|>\|f\|_\infty+\varepsilon$), so $C[0,1]$ is a
subspace of $L^\infty$ with the same norm, and $g(f)=f(0)$ has norm $1$ on it. Corollary~\ref{cor:hbnorm}(a) extends it to $\Lambda$ of norm $1$. If $\Lambda(f)=\int fg$ with $g\in L^1$, take $f_n$ continuous with values in
$[0,1]$, $f_n(0)=1$, $f_n=0$ on $[\frac1n,1]$. Then $f_ng\to0$ a.e. with $|f_ng|\le|g|$, so $\int f_ng\to0$ by dominated convergence, while $\Lambda(f_n)=1$. Since $(L^1)^*=L^\infty$ (Theorem~\ref{thm:lpdual}), the elements of $J(L^1)\subseteq(L^1)^{**}=(L^\infty)^*$ are exactly the forms $f\mapsto\int fg$, and $\Lambda$ is not one of
them.
\end{solution}

\begin{solution}{ex:mazur}
\idea Weak convergence is invisible to the norm, but convex sets are visible to linear forms: a closed convex set is cut out by half-spaces. So if no convex combination came close, a linear form would separate $x$ from all
the $x_n$, and weak convergence would fail.
\writeup Regard $E$ as a real space; real-linear continuous forms are the real parts of elements of $E^*$, and weak convergence gives $\operatorname{Re}\varphi(x_n)\to\operatorname{Re}\varphi(x)$. Let $C$ be the closure of the convex hull
of $\{x_n\}$, a closed convex set. If $x\notin C$, Theorem~\ref{thm:separation}(b) with $B=\{x\}$ gives a real form $f$ and $\alpha<\beta$ with $f(x_n)\le\alpha$ for all $n$ and $f(x)\ge\beta$, contradicting $f(x_n)\to f(x)$. So $x\in C$.
In $\ell^2$, $\frac1n(e_1+\cdots+e_n)$ has norm $\frac1{\sqrt n}\to0$: averaging recovers the weak limit $0$ in norm.
\end{solution}

\begin{solution}{ex:reflexiveweak}
\idea Compactness lives in duals (Banach--Alaoglu). A reflexive space is a dual, of its own dual, and weak* convergence there is weak convergence in $E$.
\writeup Let $\|x_n\|\le M$. Then $\|J(x_n)\|\le M$ in $E^{**}=(E^*)^*$, and $E^*$ is separable, so by Theorem~\ref{thm:alaoglu} a subsequence $J(x_{n_k})$ converges weak* to some $\Phi\in E^{**}$. By reflexivity $\Phi=J(x)$, and
$J(x_{n_k})\overset*\rightharpoonup J(x)$ means $\varphi(x_{n_k})\to\varphi(x)$ for all $\varphi\in E^*$: $x_{n_k}\rightharpoonup x$. For $E=L^p(\R^d)$, $E$ is reflexive (Proposition~\ref{prop:reflexive}) and $E^*=L^q(\R^d)$ with $q<\infty$ is separable:
step functions on boxes with rational corners and rational values are dense (Volume~6, Proposition~9.6 and Theorem~9.7). The hypothesis that $E^*$ is separable can be removed (Eberlein--\v Smulian), but not reflexivity:
in $\ell^1$, the bounded sequence $e_n$ has no weakly convergent subsequence.
\end{solution}

\begin{solution}{ex:l1quotient}
\idea A dense set of the ball is a ``basis'' in a weak sense: every vector is a fast-converging series of its elements, with coefficients halving at each step. The coefficients are a sequence in $\ell^1$.
\writeup $\|\sum_na_nx_n\|\le\sum_n|a_n|$, so $T$ is well defined, linear and $\|T\|\le1$. Let $\|y\|\le1$. Choose $n_1$ with $\|y-x_{n_1}\|<\frac12$; then $2(y-x_{n_1})$ is in the unit ball, so choose $n_2$ with
$\|2(y-x_{n_1})-x_{n_2}\|<\frac12$, that is $\|y-x_{n_1}-\frac12x_{n_2}\|<\frac14$; inductively, $\|y-\sum_{k\le m}2^{1-k}x_{n_k}\|<2^{-m}$. The sequence $a=\sum_k2^{1-k}e_{n_k}$ (indices may repeat; coefficients add) has
$\|a\|_1\le\sum_k2^{1-k}=2$, and $Ta=y$ by continuity of $T$. Hence $T$ maps the ball of radius $2$ onto a set containing the unit ball of $E$: $T$ is onto, and open. The induced map $\ell^1/\ker T\to E$ is a bounded
linear bijection, with bounded inverse by the same estimate: $\|y\|\le1$ has a preimage of norm at most $2$.
\end{solution}

\section*{Sources and further reading}

\begin{itemize}
\item S.~Banach, \emph{Th\'eorie des op\'erations lin\'eaires}, Monografje Matematyczne, Warsaw, 1932. The founding book.
\item H.~Hahn, ``\"Uber lineare Gleichungssysteme in linearen R\"aumen'', \emph{J. reine angew. Math.} 157 (1927), 214--229; S.~Banach, ``Sur les fonctionnelles lin\'eaires'', \emph{Studia Math.} 1 (1929), 211--216 and 223--239.
\item S.~Banach and H.~Steinhaus, ``Sur le principe de la condensation de singularit\'es'', \emph{Fund. Math.} 9 (1927), 50--61.
\item R.~D.~Mauldin (ed.), \emph{The Scottish Book: Mathematics from the Scottish Caf\'e}, 2nd ed., Birkh\"auser, 2015. Problem~153 and its goose are there.
\item P.~Enflo, ``A counterexample to the approximation problem in Banach spaces'', \emph{Acta Math.} 130 (1973), 309--317.
\item H.~Brezis, \emph{Functional Analysis, Sobolev Spaces and Partial Differential Equations}, Springer, 2011, Chapters~1--4. Our proofs of the separation and open mapping theorems follow it.
\item W.~Rudin, \emph{Functional Analysis}, 2nd ed., McGraw--Hill, 1991, Chapters~2--4; and \emph{Real and Complex Analysis}, 3rd ed., 1987, Chapter~6 for the duals of $L^p$ and $C(K)$.
\item L.~Carleson, ``On convergence and growth of partial sums of Fourier series'', \emph{Acta Math.} 116 (1966), 135--157.
\item S.~M.~Ulam, \emph{Adventures of a Mathematician}, Scribner, 1976, for the sessions at the Scottish Caf\'e (one of them lasted seventeen hours) and
the origin of the \emph{Scottish Book}; the Enflo goose ceremony is described in Mauldin's edition, cited above.
\end{itemize}
